\appendix

\section{Depth 21 Under the Single $K$-Cycle}
\label{app:period}

Under a single 8-cycle, $\pi^{21} = \pi^{21 \bmod 8} = \pi^{5}$.
Depth 21 therefore shares its target with depth 5; the two are not
distinct group-theoretic queries. The depth-21 probe measures
stability under 21 sequential applications of the trained operator.
In the four converged seeds, $\recninety \geq 0.9996$ at both depths,
so the additional applications preserve the high recovered fraction.
% <!-- evidence: R3 -->
For the rank-capped operator, the recovered fraction falls from
0.947 at $h{=}5$ to 0.060 at $h{=}21$.
% <!-- evidence: R4 -->
This comparison tests error accumulation, rather than generalization
to an unseen permutation target.

\begin{table}[H]
\centering
\small
\begin{tabular}{cccc}
\toprule
$h$ & predicted $\cos$ & measured $\overline{\cos}$ & measured $\recninety$ \\
\midrule
1  & 0.9317 & 0.9303 & 0.996 \\
3  & 0.9258 & 0.9259 & 0.986 \\
5  & 0.9181 & 0.9212 & 0.947 \\
7  & 0.9089 & 0.9163 & 0.881 \\
21 & 0.7536 & 0.8206 & 0.060 \\
\bottomrule
\end{tabular}
\caption{Recovery across depths for the operator with $d=16$, $K=8$,
and rank cap $k = 7 = K - 1$. Cosine predicted using the entity-subspace operator and ideal cycle
(without raw keys) is compared with measured mean cosine and the
fraction of queries with cosine greater than 0.9. Prediction error
is below 0.008 through $h = 7$. At $h = 21$, mean cosine remains
0.8206, but only 0.060 of queries exceed the recovery threshold.}
% <!-- evidence: R4 -->
\label{tab:depthcurve}
\end{table}

\Needspace{24\baselineskip}
\section{Per-Seed Subspace Decomposition}
\label{app:subspace}

\begin{table}[H]
\centering
\small
\begin{tabular}{lccccc}
\toprule
seed & $\mathrm{effrank}(A)$ & residual & leakage (\% of $\|Z\|$) & $\mathrm{effrank}(D)$ & $\rho(D)$ \\
\midrule
s1 & 8.000 & 0.7\% & 0.44\% & 8.000 & 1.38 \\
s2 & 7.999 & 1.9\% & 1.06\% & 8.000 & 1.27 \\
s3 & 7.999 & 1.5\% & 0.87\% & 8.000 & 2.86 \\
s4 & 7.999 & 2.4\% & 1.45\% & 7.999 & 1.02 \\
\bottomrule
\end{tabular}
\caption{Entity-subspace decomposition of four converged $d{=}16$, $K{=}8$
composition seeds, recomputed from archived matrices and averaged
over four evaluation episodes. $A$ acts within the entity subspace.
Residual is the scale-corrected Frobenius distance to the ideal cycle,
after fitting the isotropic scale that cosine scoring does not
measure. Leakage is the cross-block norm $\|C\|$ as a fraction of
$\|Z\|$. $D$ acts within the complementary subspace;
$\rho(D)$ is its spectral radius. The complement is full-rank and
non-contractive, while the measured cross-block leakage is small
but nonzero.}
% <!-- evidence: R5 -->
\label{tab:subspace}
\end{table}

\section{Reproducibility}
\label{app:repro}

Training, evaluation, and analysis code and archived results are
available in the public repository.\footnote{\url{https://github.com/saml212/matrix-memories}}
Per-run JSON archives support the reported numbers. A versioned script regenerates each figure
and verifies the checksum of every source file it loads. Hypotheses,
metrics, rank controls, and decision bands were registered in design
documents before training. One metric was re-registered before the
affected sweep; that deviation is recorded in those documents.

The pre-registered criteria for challenging the empirical conclusions
remain open: a rank cap of $K{-}1$ achieving at least
$0.9\times$ the converged unconstrained model's recovery at multiple
$(d, K)$ points, or a $d \geq 32$ configuration clearing the 0.7 recovery
threshold. These criteria concern measured recovery under the
experimental conditions, not the algebraic requirement for exact
linear recovery.

The rank-cap sweep aggregate retains one recovery value per
$(d, K, k)$ configuration, without per-seed instability diagnostics.
It therefore describes the observed recovery transition but cannot
separate representational limits from numerical training failures
in individual runs. Variation above the transition should likewise
be interpreted at the aggregate level.

On some platforms, the analysis pipeline emits BLAS
\texttt{RuntimeWarning}s for near-singular intermediate products.
Every decomposition output is checked for finiteness before use.
